\documentclass[11pt,a4paper]{article}
\usepackage[margin=25mm]{geometry}
\usepackage{fontspec}
\setmainfont{TeX Gyre Pagella}
\usepackage{amsmath,amssymb,amsthm,mathtools}
\usepackage{enumitem,microtype,xcolor}
\usepackage[colorlinks=true,linkcolor=blue!45!black,citecolor=blue!45!black,urlcolor=blue!45!black]{hyperref}
\usepackage{fancyhdr}
\pagestyle{fancy}
\fancyhf{}
\fancyhead[L]{\small Boundary and nonboundary localization quotients}
\fancyhead[R]{\small Unified research note}
\fancyfoot[C]{\thepage}
\setlength{\headheight}{14pt}
\setlength{\emergencystretch}{3em}
\setlength{\parskip}{3pt}
\setlist{itemsep=3pt,topsep=5pt}
\allowdisplaybreaks[2]
\newcommand{\C}{\mathbb C}
\newcommand{\Z}{\mathbb Z}
\newcommand{\N}{\mathbb N}
\newcommand{\g}{\mathfrak g}
\newcommand{\slc}{\mathfrak{sl}_2}
\newcommand{\slh}{\widehat{\mathfrak{sl}}_2}
\newcommand{\slt}{\widetilde{\mathfrak{sl}}_2}
\newcommand{\ad}{\operatorname{ad}}
\newcommand{\Ind}{\operatorname{Ind}}
\newcommand{\Ann}{\operatorname{Ann}}
\newcommand{\Hom}{\operatorname{Hom}}
\newcommand{\Span}{\operatorname{span}}
\newcommand{\Supp}{\operatorname{Supp}}
\newcommand{\pole}{\operatorname{pole}}
\newcommand{\I}{\mathcal I}
\newcommand{\X}{\mathcal X}
\newcommand{\D}{\mathcal D}
\newcommand{\T}{\mathcal T}
\newtheoremstyle{plainupright}{6pt}{5pt}{\normalfont}{}{\bfseries}{.}{.5em}{}
\theoremstyle{plainupright}
\newtheorem{proposition}{Proposition}[section]
\newtheorem{lemma}[proposition]{Lemma}
\newtheorem{theorem}[proposition]{Theorem}
\newtheorem{corollary}[proposition]{Corollary}
\newtheorem{remark}[proposition]{Remark}
\hypersetup{pdftitle={Boundary and Nonboundary Localization Quotients},
 pdfsubject={Unified note: boundary results of living note v8 and nonboundary results of 2026-09-11}}
\title{\vspace{-14mm}Boundary and Nonboundary Localization Quotients\\
of Smooth Affine \(\slh\)-Modules\\[3pt]
\large Unified note: the boundary, simplicity, and flow results of the living note v8 (2026-09-02)\\
together with the nonboundary results of \emph{nonboundary\_localization\_results} (2026-09-11)}
\author{}
\date{September 2026}
\begin{document}
\maketitle
\vspace{-9mm}
\begin{abstract}
This note consolidates the two current strands of the project. The boundary
strand (living note v8, \cite{LivingV8}) proves
\(\T_{f_b}M^c_{a,b}(\varphi)\cong M^c_{a+1,b-1}(\psi)\), its simplicity flow,
the derivation version, the iterated boundary chains, the \(\Omega\)-quotient
flow, and the critical-level flow with preserved \(\theta\)-parameters. The
nonboundary strand (\cite{Nonboundary}) studies the quotients
\(Q_c=\T_{f_c}M_{a,b}\) with \(c<b\): PBW models, root-mode and Cartan
profiles, weight structure, cyclicity, an explicit proper submodule at
integral parameters in the deep region, a Takiff model of the basic module,
annihilator preservation, and derived-functor properties. The two strands are
joined by the trichotomy for arbitrary modes, the commuting-localization
lemma, and the fact that iterating boundary moves fills any interval of
missing modes. All statements have direct proofs; this note is not a claim of
independent peer review or formal verification.
\end{abstract}

\section{Setup and background}\label{sec:setup}

All vector spaces and tensor products are over \(\C\). Let
\[
\slh=(\slc\otimes\C[t,t^{-1}])\oplus\C K,\qquad
\slt=\slh\oplus\C d,\qquad U=U(\slh),
\]
with \(x_r=x\otimes t^r\) and the normalization
\begin{align*}
[e_r,f_s]&=h_{r+s}+r\delta_{r+s,0}K,&
[h_r,h_s]&=2r\delta_{r+s,0}K,\\
[h_r,e_s]&=2e_{r+s},&
[h_r,f_s]&=-2f_{r+s},&
[d,x_r]&=rx_r.
\end{align*}
The \(e\)'s commute with one another, the \(f\)'s commute with one another,
and \(K\) is central. Fix \(a,b\in\Z\) with \(a+b\ge0\), set \(L=a+b+1\ge1\),
and put
\[
S_{a,b}=\Span\{K,\ h_i\ (i\ge0),\ e_i\ (i>a),\ f_j\ (j>b)\}.
\]
Fix a character \(\chi:S_{a,b}\to\C\) with \(\chi(e_i)=\chi(f_j)=0\),
\(\chi(K)=\kappa\), \(\chi(h_i)=\lambda_i\), and \(\lambda_i=0\) for \(i>L\).
These conditions make \(\chi\) a Lie algebra character. Let
\[
M=M_{a,b}(\chi)=U\otimes_{U(S_{a,b})}\C_\chi,\qquad
M^f_{a,b}(\chi)=\Ind_{S_{a,b}}^{\slt}\C_\chi,\qquad u=1\otimes1_\chi.
\]
For an injective root vector \(F\) on a module \(V\), set
\[
D_FV=U_F\otimes_UV,\qquad U_F=U[F^{-1}],\qquad
\T_FV=D_FV/V.
\]
More generally, for an arbitrary \(U\)-module \(V\), define
\(\T_FV\) as the cokernel of the localization map \(V\to D_FV\); then
\(\T_FV=(U_F/U)\otimes_UV\).

Fix an integer \(c<b\), write
\[
F=f_c,\qquad d_*=b-c\ge1,\qquad
Q_c=\T_FM=D_FM/M,\qquad w_m=F^{-m}u+M\ (m\ge1).
\]
The powers of \(F\) form an Ore set because \(\ad F\) is locally nilpotent on
\(U\). \emph{There is no operator \(F^{-1}\) on \(Q_c\).}

A module is \emph{smooth} if each vector is annihilated by
\(\slc\otimes t^N\C[t]\) for some \(N\). An operator \(x\) is locally
nilpotent if each vector is killed by a power of \(x\), and locally finite if
each vector spans a finite-dimensional \(\C[x]\)-module. A vector \(v\) is
\(F\)-torsion if \(F^nv=0\) for some \(n\). Weights below refer to
\(\C h_0\oplus\C K\).

The original induced modules satisfy the FGXZ criterion \cite{FGXZ}:

\begin{theorem}[FGXZ criterion]\label{thm:fgxz}
\(M_{a,b}(\chi)\) is simple if and only if
\[
\lambda_L\ne0\qquad\text{and}\qquad \kappa\ne-2.
\]
\end{theorem}

This criterion cannot be transferred directly to \(Q_c\).

\section{PBW bases and localization normal forms}\label{sec:pbw}

\begin{proposition}[Left-normal PBW model and pole filtration]\label{prop:pbw}
Choose a PBW order with \(F=f_c\) first, followed by the free families
\[
f_j\ (j\le b,\ j\ne c),\qquad e_i\ (i\le a),\qquad h_r\ (r<0),
\]
and finally a basis of \(S_{a,b}\). Let \(\X\) denote the ordered monomials in
the three middle families, including the empty monomial. Then
\begin{align*}
M &: \quad \{F^qXu:q\ge0,\ X\in\X\},\\
D_FM &: \quad \{F^qXu:q\in\Z,\ X\in\X\},\\
Q_c &: \quad \{F^{-m}Xu+M:m\ge1,\ X\in\X\}
\end{align*}
are bases. In particular, \(Q_c\) is nonzero and countably
infinite-dimensional. The operator \(F\) on \(Q_c\) is surjective and locally
nilpotent, and
\[
\ker F^r=\Span\{F^{-m}Xu+M:1\le m\le r,\ X\in\X\}\qquad(r\ge1).
\]
Every nonzero \(\g\)-submodule of \(Q_c\) meets \(\ker F\) nontrivially.
\end{proposition}

\begin{proof}
The first basis is the induced-module PBW theorem, and it makes \(M\) free
over \(\C[F]\). Localization gives the second basis: spanning follows by
moving denominators to the left, and linear independence by multiplying a
putative relation by a common power of \(F\). Quotienting gives the third.
The kernel formulas are then immediate; for the last assertion, apply to a
nonzero vector of a submodule the highest power of \(F\) that does not kill
it.
\end{proof}

\begin{lemma}[Commutation and generation by pure fractions]\label{lem:fractions}
For \(s\in U\) and \(m\ge1\),
\begin{equation}\label{eq:commute}
sF^{-m}=\sum_{j\ge0}\binom{m+j-1}{j}F^{-m-j}(\ad F)^j(s),
\end{equation}
a finite sum. Moreover, \(Q_c=\sum_{m\ge1}Uw_m\).
\end{lemma}

\begin{proof}
For \(m=1\), commute \(s\) past \(F^{-1}\) repeatedly; the process terminates
because \(\ad F\) is locally nilpotent, and induction on \(m\) gives the
binomial coefficient. For the generation assertion, give a numerator PBW word
the degree \(\rho=2\#e+\#h\), with \(f\)-factors of degree zero. Reordering
does not increase this degree, and every nonzero application of \(\ad F\)
lowers it strictly. Thus
\[
Xw_m=F^{-m}Xu+M+\text{a finite sum of terms of strictly smaller numerator degree}.
\]
Induct on \(\rho(X)\), simultaneously for all denominators \(m\).
\end{proof}

\begin{lemma}[Right-normal localized PBW]\label{lem:right-pbw}
Let \(F=f_r\) with \(r\le b\). Choose a PBW order for \(M^c_{a,b}(\chi)\) in
which \(f_r\) is the \emph{last} free variable, and let \(\X_r\) denote the
remaining ordered monomials. Then
\[
\{XF^qu:X\in\X_r,\ q\in\Z_{\ge0}\},\qquad
\{XF^qu:X\in\X_r,\ q\in\Z\},\qquad
\{XF^{-m}u+M:X\in\X_r,\ m\ge1\}
\]
are bases of \(M^c_{a,b}(\chi)\), \(D_FM^c_{a,b}(\chi)\), and
\(\T_FM^c_{a,b}(\chi)\), respectively. Correspondingly,
\(\{d^kXF^qu:k\ge0,\ X\in\X_r,\ q\in\Z\}\) is a basis of
\(D_FM^f_{a,b}(\chi)\), and in the quotient only \(q<0\) remains.
\end{lemma}

\begin{proof}
\(\ad F\) is locally nilpotent on \(U\) and on \(U(\slt)\), and
\([F,d]=-rF\), so \(\{F^m:m\ge0\}\) is an Ore set. The first basis is PBW with
\(f_r\) last. In the localization, the finite commutation formulas move
negative powers to the right of every ordered monomial. Linear independence
follows on the associated graded, where the localization is the localization
of a polynomial ring at \(\operatorname{gr}F\). Placing \(d\) first gives the
last assertions.
\end{proof}

\section{The boundary case}\label{sec:boundary}

Set
\[
F=f_b,\qquad E=e_{a+1},\qquad H=h_{a+b+1}=h_L.
\]
Define \(\psi:S_{a+1,b-1}\to\C\) by
\[
\psi(K)=\chi(K),\qquad
\psi(h_0)=\chi(h_0)+2,\qquad
\psi(h_i)=\chi(h_i)\ (i>0),
\]
and zero on the root vectors of \(S_{a+1,b-1}\). Thus \(\psi\) differs from
\(\chi\) only in the \(h_0\)-eigenvalue.

\begin{lemma}[\(\psi\) is a character]\label{lem:char}
The displayed assignment defines a Lie algebra character
\(\psi:S_{a+1,b-1}\to\C\).
\end{lemma}

\begin{proof}
The only bracket that can fail is
\([e_i,f_j]=h_{i+j}+i\delta_{i+j,0}K\) with \(i>a+1\), \(j>b-1\). Then
\(i+j\ge a+b+2=L+1\), so the central term vanishes and
\(\psi(h_{i+j})=\lambda_{i+j}=0\). The brackets \([h_i,e_j]\), \([h_i,f_j]\)
land in annihilated tails, and \([h_i,h_j]=2i\delta_{i+j,0}K\) with
\(i,j\ge0\) has value zero.
\end{proof}

\begin{lemma}[Cyclic vector]\label{lem:cyc}
The vector \(\bar u=F^{-1}u+M\in\T_FM\) satisfies
\[
x\bar u=\psi(x)\bar u\qquad(x\in S_{a+1,b-1}).
\]
\end{lemma}

\begin{proof}
For \(j>b-1\): if \(j=b\), then \(f_b\bar u=u+M=0\); if \(j>b\), then
\(f_ju=0\). From \(h_iF^{-1}=F^{-1}h_i+2F^{-2}f_{i+b}\) one gets
\(h_0\bar u=(\lambda_0+2)\bar u\), since \(2F^{-2}f_bu=2F^{-2}Fu=2\bar u\),
and \(h_i\bar u=\lambda_i\bar u\) for \(i>0\). For \(j>a+1\),
\[
e_jF^{-1}=F^{-1}e_j-F^{-2}h_{j+b}-2F^{-3}f_{j+2b},
\]
with \(j+b>L\), so \(h_{j+b}u=0\), and \(j+2b>b\), so \(f_{j+2b}u=0\). Hence
\(e_j\bar u=0\), and \(K\bar u=\kappa\bar u\).
\end{proof}

\begin{theorem}[Boundary localization]\label{thm:boundary}
If
\[
\lambda_L=\chi(h_L)\ne0,
\]
then there is a natural \(\slh\)-module isomorphism
\[
\boxed{\ \T_{f_b}M^c_{a,b}(\chi)\cong M^c_{a+1,b-1}(\psi),\ }
\]
sending the standard generator \(u'\) of the target to
\[
\bar u=f_b^{-1}u+M.
\]
The statement is independent of the level \(\kappa\).
\end{theorem}

\begin{proof}
By Lemma~\ref{lem:cyc} and the universal property of induction, there is a
unique \(\slh\)-module map
\[
\Phi:M^c_{a+1,b-1}(\psi)\longrightarrow\T_FM,
\qquad u'\longmapsto\bar u.
\]
Let \(\X\) be the ordered monomials in \(e_i\ (i\le a)\), \(h_j\ (j<0)\),
\(f_k\ (k\le b-1)\), and choose the source PBW order with \(E\) last. Then
\(\{XE^pu':X\in\X,\ p\ge0\}\) is a basis of \(M^c_{a+1,b-1}(\psi)\), and by
Lemma~\ref{lem:right-pbw}
\(\{XF^{-p-1}u+M:X\in\X,\ p\ge0\}\) is a basis of the quotient. The relations
\[
[E,F]=H,\qquad (\ad F)^2(E)=-2f_{a+2b+1},
\]
together with \(Eu=0\), \(Hu=\lambda_Lu\), and \(f_{a+2b+1}u=0\), give
\[
EF^su=s\lambda_LF^{s-1}u\qquad(s\in\Z),
\]
and in particular
\[
E^pF^{-1}u=(-1)^pp!\lambda_L^{\,p}F^{-p-1}u.
\]
Therefore
\[
\Phi(XE^pu')=(-1)^pp!\lambda_L^{\,p}XF^{-p-1}u+M,
\]
so \(\Phi\) sends a PBW basis bijectively onto nonzero scalar multiples of a
PBW basis and is an isomorphism.
\end{proof}

\begin{remark}[Surjection plus simplicity shortcut]\label{rem:shortcut}
If only simplicity is needed, one may avoid the base comparison. The chain
\(EF^{-m}u=-m\lambda_LF^{-m-1}u\) gives all pure fractions in the image of
\(\Phi\), and Lemma~\ref{lem:fractions} gives surjectivity. If \(\kappa\ne-2\),
then by Theorem~\ref{thm:fgxz} the source \(M^c_{a+1,b-1}(\psi)\) is simple,
because \(\psi(h_L)=\lambda_L\ne0\) and \(\psi(K)=\kappa\ne-2\). A nonzero
surjection from a simple module is an isomorphism.
\end{remark}

\begin{corollary}[Simplicity of the boundary quotient]\label{cor:simple}
If \(\lambda_L\ne0\) and \(\kappa\ne-2\), then \(\T_{f_b}M^c_{a,b}(\chi)\) is
simple: it is isomorphic to \(M^c_{a+1,b-1}(\psi)\), whose two irreducibility
parameters are
\[
\psi(h_{a+b+1})=\lambda_L,\qquad \psi(K)=\kappa.
\]
\end{corollary}

\begin{remark}[Why not the full localization]\label{rem:full}
The full localization \(D_FM\) contains \(M\) as a nonzero submodule, so it
cannot be simple. The object carrying the standard-module structure is the
quotient \(\T_FM=D_FM/M\). Twisted localizations keep \(F\) invertible and
are likewise not standard induced modules.
\end{remark}

\begin{proposition}[Boundary \(b\)-function]\label{prop:bfun}
With \(F=f_b\), \(E=e_{a+1}\),
\[
EF^su=s\chi(h_L)F^{s-1}u\qquad(s\in\Z),
\]
equivalently \(EF^{s+1}u=(s+1)\chi(h_L)F^su\), so the monic \(b\)-polynomial is
\[
b(s)=s+1,\qquad\text{or, with parameters,}\qquad
b_\chi(s)=(s+1)\chi(h_L).
\]
The condition \(\lambda_L\ne0\) is exactly the absence of obstruction in the
negative-power direction.
\end{proposition}

\begin{proof}
This is the finite expansion
\(EF^s=F^sE+sF^{s-1}H-2\tbinom{s}{2}F^{s-2}f_{a+2b+1}\) evaluated on \(u\).
See \cite[Section 5]{FGM} for the analogous resonance for free-field modules.
\end{proof}

\section{Arbitrary modes: trichotomy and the gap invariant}\label{sec:trichotomy}

\begin{proposition}[Arbitrary-mode trichotomy]\label{prop:trichotomy}
Let \(M=M^c_{a,b}(\chi)\) and consider \(F=f_r\).
\begin{enumerate}[label=\textup{(\roman*)},leftmargin=2.8em]
\item If \(r>b\), then \(F\) is locally nilpotent on \(M\), hence
\(D_FM=0\): there is no nontrivial localization.
\item If \(r=b\) and \(\chi(h_L)\ne0\), then
\(\T_{f_b}M\cong M^c_{a+1,b-1}(\psi)\) by Theorem~\ref{thm:boundary}.
\item If \(r<b\), then \(D_FM\ne0\), but \(\T_FM\) is not isomorphic to any
standard module \(M^c_{a',b'}(\chi')\): the set of locally nilpotent negative
root modes is \(\{r\}\cup\{j:j>b\}\), which has a gap, while standard modules
have full upper tails \(\{j:j>b'\}\).
\end{enumerate}
The same trichotomy holds for \(M^f_{a,b}(\chi)\). Case (iii) is developed in
depth in Sections~\ref{sec:cartan}--\ref{sec:deep}.
\end{proposition}

\begin{proof}
For \(r>b\), \(Fu=0\) and \(F^m(xu)=(\ad F)^m(x)u\) vanishes for \(m\gg0\),
so \(F\) is locally nilpotent on \(M\). For \(r<b\), Lemma~\ref{lem:right-pbw}
gives the localized basis; \(F\) is locally nilpotent on the quotient, \(f_j\)
for \(j\le b\), \(j\ne r\) acts as an injective free polynomial variable, and
\(f_j\) for \(j>b\) inherits local nilpotence from \(M\).
\end{proof}

\begin{proposition}[The negative-root profile]\label{prop:root}
On \(Q_c\),
\[
\begin{array}{c|l}
j=c&f_j\text{ is surjective and locally nilpotent},\\
j\le b,\ j\ne c&f_j\text{ is injective and not surjective},\\
j>b&f_j\text{ is locally nilpotent}.
\end{array}
\]
Consequently the set of locally nilpotent negative-root modes is exactly
\(\{c\}\cup\{j:j>b\}\).
\end{proposition}

\begin{proof}
The first assertion is Proposition~\ref{prop:pbw}. For \(j\le b\), \(j\ne c\),
\(f_j\) is multiplication by a free polynomial variable in the commuting
algebra model. For \(j>b\), use \(f_ju=0\), local nilpotence of \(\ad f_j\),
and \([f_j,F]=0\).
\end{proof}

\begin{corollary}[Nonisomorphism and Hom vanishing]\label{cor:hom}
For \(c<b\), \(Q_c\) is not isomorphic to any standard module
\(M_{a',b'}(\psi)\). For fixed \(M\) and distinct \(c,c'\le b\),
\[
\Hom_\g(T_{f_c}M,T_{f_{c'}}M)=0.
\]
\end{corollary}

\begin{proof}
The local-nilpotence set of a standard module is a half-line. For the Hom
assertion, \(f_c\) is locally nilpotent on the source and injective on the
target.
\end{proof}

\begin{proposition}[Smoothness and weight structure]\label{prop:smooth}
The module \(Q_c\) is smooth, \(K\) acts by \(\kappa\), and \(h_0\) acts
semisimply with
\[
\Supp_{h_0}Q_c=\lambda_0+2\Z.
\]
Every nonzero weight space is countably infinite-dimensional. The module has
no nonzero finite-dimensional submodules or quotients, and is not integrable
for the full affine algebra.
\end{proposition}

\begin{proof}
If \(X\) has \(r_e\) \(e\)-factors and \(r_f\) \(f\)-factors, then
\[
h_0(F^{-m}Xu+M)=(\lambda_0+2m+2r_e-2r_f)(F^{-m}Xu+M).
\]
All weights in the stated coset occur; appending powers of \(h_{-1}\) gives
infinitely many independent vectors of each weight. Smoothness follows from
the finite commutation formulas
\[
h_rF^{-m}=F^{-m}h_r+2mF^{-m-1}f_{r+c},
\]
\[
e_rF^{-m}=F^{-m}e_r-mF^{-m-1}(h_{r+c}+r\delta_{r+c,0}K)
-m(m+1)F^{-m-2}f_{r+2c},
\]
by taking \(r,r+c,r+2c\) sufficiently large. Surjectivity and local
nilpotence of \(F\) exclude nonzero finite-dimensional quotients, and a
finite-dimensional smooth \(\g\)-module is trivial; trivial submodules are
excluded because \(f_{c+1}\) is injective. A free root mode \(f_j\),
\(j\le b\), \(j\ne c\), excludes integrability.
\end{proof}

\section{Positive Cartan modes and localization depth}\label{sec:cartan}

\begin{proposition}[Strict growth of pole order]\label{prop:pole}
For \(1\le i\le d_*\), \(\zeta\in\C\), and \(0\ne v\in Q_c\),
\[
\pole_F((h_i-\zeta)v)=\pole_F(v)+1,
\]
where pole order means the largest denominator exponent in the PBW expansion.
In particular,
\[
\ker p(h_i)=0\qquad(0\ne p\in\C[t]).
\]
Thus \(h_i\) has no nonzero locally finite vectors, eigenvectors, or
generalized eigenvectors.
\end{proposition}

\begin{proof}
Write the top pole term of \(v\) as \(F^{-m}v_m+M\), \(v_m\ne0\).
The term \(F^{-m}h_iv_m\) has pole order at most \(m\), and the coefficient of
pole order \(m+1\) is \(2mf_{c+i}v_m\pmod{FM}\). Since \(c+i\le b\) and
\(c+i\ne c\), multiplication by \(f_{c+i}\) is injective on \(M/FM\), so this
coefficient is nonzero and cannot be cancelled by lower terms or by
\(-\zeta v\). Iteration shows that \(v,h_iv,\ldots,h_i^nv\) have distinct pole
orders.
\end{proof}

\begin{corollary}\label{cor:no-whittaker}
For every standard induced module,
\[
\Hom_\g(M_{a',b'}(\psi),Q_c)=0.
\]
More generally, \(Q_c\) has no nonzero character vector for any Lie
subalgebra containing \(h_1\).
\end{corollary}

\begin{proposition}[High Cartan modes]\label{prop:high}
If \(i>d_*\), then \(h_iw_m=\lambda_iw_m\) for every \(m\ge1\), and
\(h_i-\lambda_i\) is locally nilpotent on all of \(Q_c\).
\end{proposition}

\begin{proof}
Since \(c+i>b\), the correction term in the \(h\)-commutation kills \(u\).
For a general smooth module, if \(i>0\) and \(h_iw=\lambda w\), then
\((h_i-\lambda)^NXw=(\ad h_i)^N(X)w\); for large \(N\) every surviving term
contains a root factor of arbitrarily high index, which annihilates \(w\).
Apply this to all \(w_m\) and use Lemma~\ref{lem:fractions}.
\end{proof}

Consequently the positive Cartan modes with no nonzero locally finite vectors
are precisely \(h_1,\ldots,h_{b-c}\): a second intrinsic way to read off the
localization depth.

\section{The basic module \texorpdfstring{\(Q_0=T_{f_0}M_{-1,1}\)}{Q0}}\label{sec:basic}

Here \(a=-1\), \(b=1\), \(c=0\), and \(L=d_*=1\). All statements allow
arbitrary \(\lambda_0,\lambda_1,\kappa\). The pure fractions satisfy
\begin{align}
f_0w_1&=0,\qquad f_0w_m=w_{m-1}\quad(m>1),\notag\\
h_0w_m&=(\lambda_0+2m)w_m,\label{eq:basic}\\
e_0w_m&=-m(\lambda_0+m+1)w_{m+1},\label{eq:basic-e}\\
h_1w_m&=\lambda_1w_m+2mf_1w_{m+1},\notag\\
e_1w_m&=-m\lambda_1w_{m+1}-m(m+1)f_1w_{m+2}.\notag
\end{align}
Unlike boundary localization at \(f_1\), the pure-fraction recursion is
controlled by \(\lambda_0\), not by \(\lambda_1\).

\begin{proposition}[Exact cyclicity statement]\label{prop:cyclic}
The module \(Q_0\) is cyclic for all parameters, and
\[
Q_0=Uw_1\quad\Longleftrightarrow\quad
\lambda_0\notin\{-2,-3,-4,\ldots\}.
\]
If \(\lambda_0=-s\) with \(s\ge2\), then \(Q_0=Uw_s\), while \(Uw_1\) is a
proper submodule.
\end{proposition}

\begin{proof}
If \(\lambda_0\) is not a negative integer at most \(-2\), every coefficient
in \eqref{eq:basic-e}, starting at \(m=1\), is nonzero, and induction gives
\(w_m\in Uw_1\) for all \(m\); Lemma~\ref{lem:fractions} then gives
\(Q_0=Uw_1\). If \(\lambda_0=-s\), the coefficient vanishes precisely at
\(m=s-1\); starting at \(w_s\), applications of \(e_0\) produce all \(w_m\)
with \(m\ge s\), and applications of \(f_0\) produce \(w_{s-1},\ldots,w_1\),
so \(Q_0=Uw_s\). But \(e_0^{s-1}w_1=0\), so \(Uw_1\subseteq G\), while
\(w_s\notin G\), with \(G\) as in Proposition~\ref{prop:integer}.
\end{proof}

\begin{proposition}[The \(e_0\)-locally nilpotent part]\label{prop:integer}
Let
\[
G=\{v\in Q_0:e_0^Nv=0\text{ for some }N\ge1\}.
\]
This is a \(\g\)-submodule, and
\[
\begin{cases}
0\ne G\ne Q_0,&\lambda_0\in\Z,\\
G=0,&\lambda_0\notin\Z.
\end{cases}
\]
Thus \(Q_0\) is reducible for every integral \(\lambda_0\). For nonintegral
\(\lambda_0\), \(e_0\) is injective; simplicity is not asserted.
\end{proposition}

\begin{proof}
Set \(E=e_0\). Since \(\ad E\) is locally nilpotent on \(U\), the expansion
\[
E^n(sv)=\sum_{r=0}^n\binom nr(\ad E)^r(s)E^{n-r}v
\]
proves \(UG\subseteq G\).

Suppose first \(p=\lambda_0\in\Z\). Choose \(k\ge\max(1,p+2)\), and put
\(x=f_0\), \(y=f_1\), \(\eta=\lambda_1\). The subspace \(\C[x,y]u\) is stable
under \(e_0,h_0,f_0\), with
\begin{equation}\label{eq:polynomial-e}
e_0(x^jy^\ell u)
=j(p-j+1-2\ell)x^{j-1}y^\ell u+\ell\eta x^jy^{\ell-1}u.
\end{equation}
Define
\begin{equation}\label{eq:highest-vector}
v_k=\sum_{j=0}^k a_jx^jy^{k-j}u,\qquad a_0=1,\qquad
a_j=-\frac{(k-j+1)\eta}{j(p-2k+j+1)}a_{j-1}.
\end{equation}
All denominators are nonzero, since \(p-2k+j+1\le p-k+1<0\). The coefficient
of \(x^{j-1}y^{k-j}u\) in \(e_0v_k\) is
\[
j(p-2k+j+1)a_j+(k-j+1)\eta a_{j-1}=0,
\]
so \(e_0v_k=0\), and \(h_0v_k=\nu v_k\) with \(\nu=p-2k\le-2\). The rank-one
commutation formula gives
\[
e_0(f_0^{-m}v_k)=-m(\nu+m+1)f_0^{-m-1}v_k,
\]
whose coefficient vanishes at \(m=-\nu-1\). In \(Q_0\),
\[
f_0^{-1}v_k+M=f_1^kw_1\ne0,\qquad
e_0^{-\nu-1}(f_1^kw_1)=0,
\]
so \(G\ne0\). On the other hand, for \(m\ge\max(1,-p)\) every coefficient in
the successive applications of \eqref{eq:basic-e} to \(w_m\) is nonzero, so
\(w_m\notin G\): the submodule \(G\) is proper.

Now suppose \(\lambda_0\notin\Z\) and \(G\ne0\). Since \(G\) is \(h_0\)-stable
and \(h_0\) is semisimple, it contains a nonzero weight vector; applying a
power of \(e_0\) produces a nonzero \(v'\in\ker e_0\). Write \(h_0v'=\nu v'\)
and choose the least \(N\ge1\) with \(f_0^Nv'=0\). The \(\slc\)-identity gives
\[
0=e_0f_0^Nv'=N(\nu-N+1)f_0^{N-1}v',
\]
so by minimality \(\nu=N-1\in\Z\), contradicting \(\nu\in\lambda_0+2\Z\).
\end{proof}

\begin{remark}[What this submodule does and does not prove]\label{rem:G}
The submodule \(G\) is the largest submodule on which \(e_0\) is locally
nilpotent, and is locally finite for
\(\langle e_0,h_0,f_0\rangle\cong\slc\). The quotient \(Q_0/G\) is nonzero
and \(e_0\) acts injectively on it. Neither simplicity of \(G\) or \(Q_0/G\),
nor length two or indecomposability of \(Q_0\), has been proved.
\end{remark}

\begin{proposition}[An exact Takiff model]\label{prop:takiff}
The subspace
\[
V=\Span\{f_1^kw_m:m\ge1,\ k\ge0\}
\cong\C[x,x^{-1},y]/\C[x,y]
\]
is stable under the nonnegative current algebra \(\slc[t]\), with all modes of
degree at least two acting by zero; it is therefore a module for the Takiff
algebra \(\slc[t]/t^2\slc[t]\). Under \(f_1^kw_m\leftrightarrow x^{-m}y^k\),
\begin{align*}
f_0&=x,\qquad f_1=y,\\
h_0&=\lambda_0-2x\partial_x-2y\partial_y,\\
h_1&=\lambda_1-2y\partial_x,\\
e_0&=\lambda_0\partial_x+\lambda_1\partial_y
      -x\partial_x^2-2y\partial_x\partial_y,\\
e_1&=\lambda_1\partial_x-y\partial_x^2.
\end{align*}
These are exact formulas, not a degree truncation. The subspace \(V\) is not a
submodule for the full affine algebra (negative Cartan modes leave it).
\end{proposition}

\section{General nonboundary indices: deep and shallow regions}\label{sec:deep}

\subsection{Spectral-flow normalization}

For \(k\in\Z\), define the automorphism
\[
\sigma_k(e_i)=e_{i+k},\quad
\sigma_k(f_i)=f_{i-k},\quad
\sigma_k(h_i)=h_i+k\delta_{i,0}K,\quad
\sigma_k(K)=K.
\]
Our twist convention is \(x\cdot_{\rm new}v=\sigma_k(x)v\). Tracking the
inducing conditions and the localization element gives
\begin{equation}\label{eq:spectral}
(Q_c)^{\sigma_{-c}}\cong
T_{f_0}M_{a+c,b-c}(\chi'),
\qquad
\chi'(h_0)=\mu:=\lambda_0-c\kappa,
\end{equation}
with \(\chi'(K)=\kappa\) and \(\chi'(h_i)=\lambda_i\) for \(i>0\). This is an
isomorphism \emph{after twisting}, not an untwisted isomorphism. The useful
normalized coordinates are
\[
A=a+c,\qquad B=b-c=d_*,\qquad A+B=L-1,\qquad \mu=\lambda_0-c\kappa.
\]

\begin{proposition}[The deep region]\label{prop:deep}
Suppose \(c\le-a-1\), equivalently \(d_*\ge L\). Put
\[
E=e_{-c},\qquad H=h_0-cK,\qquad F=f_c.
\]
These elements form an \(\slc\)-triple, \(Eu=0\), \(Hu=\mu u\), and
\[
Ew_m=-m(\mu+m+1)w_{m+1}.
\]
The following statements hold without assuming \(\lambda_L\ne0\):
\begin{enumerate}[label=(\roman*)]
\item \(Q_c=Uw_1\) if and only if \(\mu\notin\{-2,-3,\ldots\}\).
If \(\mu=-s\), \(s\ge2\), then \(Q_c=Uw_s\ne Uw_1\).
\item The subspace
\[
G_c=\{v\in Q_c:E^Nv=0\text{ for some }N\ge1\}
\]
is a \(\g\)-submodule. It is nonzero and proper when \(\mu\in\Z\), and zero
when \(\mu\notin\Z\).
\item The operator \(E\) acts injectively on the nonzero quotient
\(Q_c/G_c\).
\end{enumerate}
\end{proposition}

\begin{proof}
The commutator \([e_{-c},f_c]=h_0-cK\) gives the rank-one triple; the
deep-region assumption is exactly \(-c>a\), so \(Eu=0\). Set \(Y_b=f_b\). On
\(\C[F,Y_b]u\),
\begin{equation}\label{eq:deep-poly}
E(F^jY_b^\ell u)
=j(\mu-j+1-2\ell)F^{j-1}Y_b^\ell u
+\ell\lambda_{d_*}F^jY_b^{\ell-1}u,
\end{equation}
which is \eqref{eq:polynomial-e} with
\((p,\eta,x,y)\) replaced by \((\mu,\lambda_{d_*},F,Y_b)\). When \(\mu\in\Z\),
the highest vector \eqref{eq:highest-vector} with these replacements gives
\(0\ne f_b^kw_1\in G_c\) with \(E^{2k-\mu-1}(f_b^kw_1)=0\); if \(d_*>L\),
then \(\lambda_{d_*}=0\) and the vector simply equals \(f_b^ku\). Pure
fractions with \(m\ge\max(1,-\mu)\) are not \(E\)-locally nilpotent, so
\(G_c\) is proper. For \(\mu\notin\Z\), the least-\(N\) argument of
Proposition~\ref{prop:integer} forces an integral highest weight, a
contradiction, so \(G_c=0\). Injectivity on \(Q_c/G_c\) follows from
\(Ev\in G_c\Rightarrow v\in G_c\).
\end{proof}

\begin{proposition}[The shallow region]\label{prop:shallow}
Suppose \(-a\le c<b\), equivalently \(1\le d_*<L\).
Then \(E=e_{-c}\) acts injectively on \(Q_c\), for every value of \(\mu\).
If \(\lambda_L\ne0\), the vector \(w_1\) generates \(Q_c\).
\end{proposition}

\begin{proof}
In this region \(e_{-c}\) belongs to the free PBW complement, so \(Eu=0\)
cannot be used. Instead set \(Y=e_{L-c}\). Since \(L-c>a\) and \(L+c>b\),
\[
Yu=0,\qquad [Y,F]=h_L,\qquad f_{L+c}u=0,
\]
and there is no central term in the middle identity because \(L\ge1\). The
\(e\)-commutation formula therefore yields
\[
Yw_m=-m\lambda_Lw_{m+1},
\]
so \(\lambda_L\ne0\) gives cyclicity via Lemma~\ref{lem:fractions}.

For injectivity of \(E\), use numerator degree \(\rho=2\#e+\#h\): on the
highest-degree part of a nonzero PBW expansion in \(Q_c\), left multiplication
by \(E\) adds a free \(e_{-c}\)-factor, increasing the degree by two, while
all commutator terms have degree at most one above the original. In the
associated graded PBW space the highest term is multiplication by a free
polynomial variable, hence injective, and it cannot cancel.
\end{proof}

\section{Extensions, annihilators, and derived functors}\label{sec:homological}

\begin{proposition}[The natural extension is nonsplit]\label{prop:extension}
The short exact sequence
\[
0\longrightarrow M\longrightarrow D_FM\longrightarrow Q_c\longrightarrow0
\]
does not split. In fact,
\[
\Hom_\g(Q_c,D_FM)=\Hom_\g(Q_c,M)=\Hom_\g(M,Q_c)=0.
\]
If \(M\) is simple, it is the unique simple submodule of \(D_FM\) and an
essential submodule; consequently \(D_FM\) is indecomposable.
\end{proposition}

\begin{proof}
The source \(Q_c\) is \(F\)-torsion, whereas \(F\) is injective on \(D_FM\)
and on \(M\); this proves the first two vanishings and excludes a section.
The third is Corollary~\ref{cor:no-whittaker}. For \(0\ne v\in D_FM\), a
sufficiently large power of \(F\) sends \(v\) to a nonzero element of \(M\).
\end{proof}

\begin{proposition}[Annihilator preservation]\label{prop:ann}
In the ordinary universal enveloping algebra,
\[
\Ann_U M=\Ann_U D_FM=\Ann_U Q_c.
\]
\end{proposition}

\begin{proof}
Let \(\Delta=\ad F\). The two-sided ideal \(\Ann_UM\) is \(\Delta\)-stable, so
\eqref{eq:commute} shows that every element of \(\Ann_UM\) annihilates
\(D_FM\); the inclusion \(M\subseteq D_FM\) gives equality of the first two
annihilators, and this common ideal is contained in \(\Ann_UQ_c\). For the
reverse inclusion, let \(s\in\Ann_UQ_c\), and choose \(t\ge0\) with
\(\Delta^{t+1}(s)=0\). For \(v\in M\) and \(m\ge1\), \(sF^{-m}v\in M\).
Multiplying by \(F^{m+t}\) gives
\[
P_v(m):=\sum_{j=0}^t\binom{m+j-1}{j}F^{t-j}\Delta^j(s)v
\ \in F^{m+t}M,
\]
a vector-valued polynomial in \(m\). Its finitely many coefficients have
bounded nonnegative \(F\)-exponents, so for all sufficiently large \(m\) the
left side lies in a space intersecting \(F^{m+t}M\) trivially; hence
\(P_v\equiv0\). Evaluating at \(m=0\) gives \(F^tsv=0\), because
\(\binom{j-1}{j}=0\) for \(j\ge1\); injectivity of \(F\) gives \(sv=0\).
\end{proof}

\begin{proposition}[Functorial and derived properties]\label{prop:derived}
For an arbitrary \(U\)-module \(V\),
\begin{align*}
T_FV&=(U_F/U)\otimes_UV,\\
L_1T_F(V)&\cong\{v\in V:F^nv=0\text{ for some }n\ge0\},\\
L_jT_F(V)&=0\qquad(j\ge2).
\end{align*}
The functor is right exact and is exact on short exact sequences whose three
terms are \(F\)-torsion-free. In particular,
\[
D_FQ_c=0,\qquad T_FQ_c=0,\qquad L_1T_F(Q_c)\cong Q_c.
\]
\end{proposition}

\begin{proof}
Ore localization is flat as a right \(U\)-module, so
\(0\to U\to U_F\to U_F/U\to0\) is a flat resolution of \(U_F/U\). Tensoring
with \(V\) identifies its degree-one homology with
\(\ker(V\to D_FV)\), the \(F\)-torsion submodule, and there is no higher
homology.
\end{proof}

\section{Further localization, the boundary chains, and iteration}\label{sec:chains}

\begin{proposition}[Commuting quotient localizations]\label{prop:commuting}
Let \(F=f_c\) and \(G=f_j\) with distinct \(c,j\le b\), and write
\(D_{F,G}M\) for simultaneous localization. There are natural isomorphisms
\[
T_G(T_FM)\cong
\frac{D_{F,G}M}{D_FM+D_GM}
\cong T_F(T_GM).
\]
The analogous statement holds for any finite family of distinct free
negative-root modes, independently of their order.
\end{proposition}

\begin{proof}
The modes commute, their localizations commute, and PBW makes \(M\) free in
the two independent polynomial variables \(F,G\); inside \(D_{F,G}M\),
\[
D_FM\cap D_GM=M.
\]
Moreover, \(G\) is injective on \(T_FM\). Localizing the exact sequence
defining \(T_FM\) at \(G\) gives
\(D_G(T_FM)\cong D_{F,G}M/D_GM\), and quotienting by the image of \(T_FM\)
gives the middle expression. Its PBW basis consists of terms with both
exponents strictly negative.
\end{proof}

\begin{corollary}[Filling the interval of missing modes]\label{cor:fill}
Assume \(\lambda_L\ne0\). Applying one quotient localization at each of
\[
f_c,f_{c+1},\ldots,f_b
\]
in any order yields
\[
M_{a+b-c+1,c-1}\bigl(\chi^{(b-c+1)}\bigr),
\qquad
\chi^{(b-c+1)}(h_0)=\lambda_0+2(b-c+1),
\]
with all other character values unchanged.
\end{corollary}

\begin{proof}
By Proposition~\ref{prop:commuting}, arrange the operations in the order
\(f_b,f_{b-1},\ldots,f_c\). At each step the next mode is a boundary mode for
the resulting standard module, and Theorem~\ref{thm:boundary} applies. The
value of \(L=a+b+1\) and the nonzero parameter \(\lambda_L\) are unchanged
throughout.
\end{proof}

For the basic example,
\[
T_{f_1}\bigl(T_{f_0}M_{-1,1}(\chi)\bigr)
\cong T_{f_0}\bigl(T_{f_1}M_{-1,1}(\chi)\bigr)
\cong M_{1,-1}(\chi^{(2)}),\qquad \lambda_1\ne0.
\]
These are iterated \emph{quotient} localizations; one must divide by the sum
of the relevant partial localizations, and simplicity of the final module does
not imply simplicity of the first quotient.

\begin{theorem}[Iterated boundary chain]\label{thm:iterated}
Assume \(\chi(h_L)\ne0\). For \(k\ge1\) define \(\chi^{(k)}\) by
\[
\chi^{(k)}(K)=\chi(K),\qquad
\chi^{(k)}(h_0)=\chi(h_0)+2k,\qquad
\chi^{(k)}(h_i)=\chi(h_i)\ (i>0).
\]
Then
\begin{equation}\label{eq:iterated}
\boxed{\ \T_{f_{b-k+1}}\cdots\T_{f_{b-1}}\T_{f_b}
M^c_{a,b}(\chi)\cong M^c_{a+k,b-k}(\chi^{(k)}),\ }
\end{equation}
with the rightmost functor acting first, using the same nondegeneracy
condition \(\chi(h_L)\ne0\) at every step. In particular, for
\(M^c_{-1,N}(\chi)\) and \(n\le N\), taking \(k=N-n+1\),
\begin{equation}\label{eq:iterated-special}
\boxed{\ \T_{f_n}\T_{f_{n+1}}\cdots\T_{f_N}
M^c_{-1,N}(\chi)\cong M^c_{N-n,n-1}(\chi^{(N-n+1)}).\ }
\end{equation}
\end{theorem}

\section{Adjoining the derivation}\label{sec:derivation}

On \(M^f_{a,b}(\chi)\cong\C[d]\otimes M^c_{a,b}(\chi)\), the \(\slt\)-action is
the semidirect one: for \(x_r\) and \(p(d)\in\C[d]\),
\[
d\bigl(p(d)\otimes v\bigr)=dp(d)\otimes v,\qquad
x_r\bigl(p(d)\otimes v\bigr)=p(d-r)\otimes x_rv.
\]
In particular, \(F^{-1}(p(d)\otimes v)=p(d+b)\otimes F^{-1}v\) for \(F=f_b\).

\begin{proposition}[Localization commutes with adjoining \(d\)]
\label{prop:loc-induction}
Let \(X\) be an \(F=f_b\)-torsion-free \(\slh\)-module. Then there are natural
\(\slt\)-module isomorphisms
\[
D_F\bigl(\Ind_{\slh}^{\slt}X\bigr)\cong\Ind_{\slh}^{\slt}(D_FX),
\qquad
\boxed{\ \T_F\bigl(\Ind_{\slh}^{\slt}X\bigr)\cong\Ind_{\slh}^{\slt}(\T_FX).\ }
\]
As vector spaces, \(\T_F(\C[d]\otimes X)\cong\C[d]\otimes\T_FX\).
\end{proposition}

\begin{proof}
Set \(A=U(\slh)\), \(B=U(\slt)\). Since \([d,F]=bF\), the derivation \(\ad d\)
extends to \(D_FA\), and PBW gives \(D_FB\cong B\otimes_AD_FA\). Hence
\(D_F(B\otimes_AX)\cong B\otimes_AD_FX\). Since \(B\) is free as a right
\(A\)-module, \(B\otimes_A-\) is exact, and applying it to
\(0\to X\to D_FX\to\T_FX\to0\) gives the quotient isomorphism.
\end{proof}

\begin{theorem}[Boundary localization with the derivation]\label{thm:boundary-f}
If \(\chi(h_L)\ne0\), then
\[
\boxed{\ \T_{f_b}M^f_{a,b}(\chi)\cong M^f_{a+1,b-1}(\psi).\ }
\]
In the PBW normal form \(M^f\cong\C[d]\otimes M^c\), the isomorphism is
termwise
\[
d^kXE^pu'\longmapsto
(-1)^pp!\lambda_L^{\,p}\,d^kXF^{-p-1}u+M^f_{a,b}(\chi).
\]
The iterated chains \eqref{eq:iterated}--\eqref{eq:iterated-special} hold
verbatim for \(M^f\).
\end{theorem}

\begin{proposition}[Nonboundary version with the derivation]\label{prop:d-nonb}
There is a natural \(\slt\)-module isomorphism
\[
\I(Q_c)\cong T_{f_c}\I(M),\qquad
\I(V)=\Ind_{\slh}^{\slt}V\cong\C[d]\otimes V,
\]
given explicitly by
\[
\alpha\bigl(F^{-m}(P(d)\otimes v)\bigr)=P(d+mc)\otimes F^{-m}v,
\]
with inverse \(P(d)\otimes F^{-m}v\mapsto F^{-m}(P(d-mc)\otimes v)\). If
\(0\ne N\ne Q_c\) is a \(\g\)-submodule, then
\(0\ne\I(N)\ne\I(Q_c)\).
\end{proposition}

\begin{proof}
On \(\I(M)\), the target of \(\I(M)\to\I(D_FM)\) has an \(F\)-invertible
action, with
\[
F(P(d)\otimes v)=P(d-c)\otimes Fv,\qquad
F^{-1}(P(d)\otimes v)=P(d+c)\otimes F^{-1}v.
\]
The universal property of localization extends the map to
\(\alpha:D_F\I(M)\to\I(D_FM)\) with the displayed formula, which respects
equivalent fractions. The algebra \(U(\slt)\) is free as a right \(U\)-module,
so \(\I\) is exact; quotienting identifies the cokernels. Exactness and
faithfulness of \(\I\) prove the last assertion.
\end{proof}

\section{Non-critical simple flow with the Casimir}\label{sec:omega}

Assume \(\kappa\ne-2\). The Casimir element \(\Omega\) of \(\slt\) commutes
with \(U(\slt)\); to simplify notation take the representative \(P(\Omega)=1\)
and write
\[
L_{a,b}(\chi;\mu):=M^f_{a,b}(\chi)/(\Omega-\mu)M^f_{a,b}(\chi).
\]

\begin{theorem}[Non-critical simple flow]\label{thm:omega-flow}
Assume
\[
\chi(h_L)\ne0,\qquad \kappa\ne-2.
\]
Then \(f_b\) is injective on \(L_{a,b}(\chi;\mu)\), and
\[
\boxed{\ \T_{f_b}L_{a,b}(\chi;\mu)\cong L_{a+1,b-1}(\psi;\mu).\ }
\]
Both sides are simple \(\slt\)-modules, and the Casimir parameter \(\mu\) does
not change. More generally,
\[
\T_{f_{b-k+1}}\cdots\T_{f_b}L_{a,b}(\chi;\mu)
\cong L_{a+k,b-k}(\chi^{(k)};\mu).
\]
\end{theorem}

\begin{proof}
By \cite[formula (4.1)]{FGXZ}, at level \(\kappa\ne-2\) the restriction of
\(L_{a,b}(\chi;\mu)\) to \(\slh\) is isomorphic to \(M^c_{a,b}(\chi)\); hence
\(f_b\) remains injective. Since \(\Omega\) commutes with \(F=f_b\), exactness
of Ore localization gives
\[
\T_F\bigl(M^f/(\Omega-\mu)M^f\bigr)
\cong
\bigl(\T_FM^f\bigr)/(\Omega-\mu)\bigl(\T_FM^f\bigr).
\]
Theorem~\ref{thm:boundary-f} turns the right side into
\(M^f_{a+1,b-1}(\psi)/(\Omega-\mu)M^f_{a+1,b-1}(\psi)
=L_{a+1,b-1}(\psi;\mu)\). The irreducibility parameters of
\cite[Theorem 4.1]{FGXZ} are unchanged, so both sides are simple. Iterate.
\end{proof}

\begin{remark}[Why one cannot stop at \(M^f\)]\label{rem:omega}
At non-critical level, \((\Omega-\mu)M^f\) is a proper submodule of \(M^f\),
so the raw module \(M^f\) and the raw quotient
\(\T_{f_b}M^f\cong M^f_{a+1,b-1}(\psi)\) are generally not simple. The
complete simple flow is
\[
\boxed{\ \text{boundary localization quotient}\ +\ \Omega\text{-central reduction},\ }
\]
and the two operations commute, with \(\mu\) preserved.
\end{remark}

\section{Critical level: the \(\theta\)-parameters are preserved}\label{sec:critical}

Throughout this section \(\kappa=\chi(K)=-2\) and
\(q:=L=a+b+1\). For the reduced case \((a,b)=(-1,q)\), the critical central
modes used in \cite{FGXZ} are the Segal--Sugawara modes \(T(q-i)\).

\begin{proposition}[Spectral flow on the Sugawara modes]\label{prop:spectral}
Let \(\sigma_k\) be as in Section~\ref{sec:deep}. In the completed enveloping
algebra,
\[
\boxed{\ \sigma_k(T(m))
=T(m)+\frac{k}{2}(K+2)h_m+\frac{k^2}{4}K(K+2)\delta_{m,0}.\ }
\]
Consequently, at the critical level \(K=-2\),
\[
\boxed{\ \sigma_k(T(m))=T(m)\ }
\]
as operators on any smooth module.
\end{proposition}

\begin{proof}
Write \(T(m)\) as the combination \(\tfrac14(2,2,1)\) of the three
normal-ordered sums of \cite[formula (2.3)]{FGXZ}. For \(k>0\), spectral flow
moves the normal-ordering cut of the first sum from \(i>0\) to \(i>-k\),
producing \(kh_m+\tfrac{k(k-1)}{2}K\delta_{m,0}\); the second sum produces
\(kh_m+\tfrac{k(k+1)}{2}K\delta_{m,0}\); the third produces
\(2kKh_m+k^2K^2\delta_{m,0}\). Inserting the coefficients \(\tfrac14(2,2,1)\)
gives the formula; negative \(k\) follows from
\(\sigma_{k+\ell}=\sigma_k\sigma_\ell\), and \(K=-2\) gives the invariance.
\end{proof}

For the critical character, define
\[
Z_i^{a,b}:=\sigma_{-a-1}\bigl(T(q-i)\bigr),\qquad i\in\N,
\]
and let \(I^{a,b}_{\boldsymbol\theta}\) be the ideal of the critical center
generated by \(Z_i^{a,b}-\theta_i\). Put
\[
\overline M^{\,c}_{a,b}(\chi;\boldsymbol\theta)
:=M^c_{a,b}(\chi)/I^{a,b}_{\boldsymbol\theta}M^c_{a,b}(\chi),
\qquad
\overline M^{\,f}_{a,b}(\chi;\boldsymbol\theta)
:=\Ind_{\slh}^{\slt}\overline M^{\,c}_{a,b}(\chi;\boldsymbol\theta).
\]
By Proposition~\ref{prop:spectral}, at level \(-2\)
\[
Z_i^{a,b}=T(q-i)=Z_i^{a+1,b-1},
\]
so the boundary move does not change \(\boldsymbol\theta\).

\begin{theorem}[Critical simple flow]\label{thm:critical-flow}
Assume
\[
\chi(K)=-2,\qquad \chi(h_q)\ne0.
\]
Then \(f_b\) is injective on \(\overline M^{\,c}_{a,b}(\chi;\boldsymbol\theta)\), and
\[
\boxed{\ \T_{f_b}\overline M^{\,c}_{a,b}(\chi;\boldsymbol\theta)
\cong \overline M^{\,c}_{a+1,b-1}(\psi;\boldsymbol\theta),\ }
\]
both sides being simple \(\slh\)-modules. The same holds with \(f\) in place
of \(c\) throughout, both sides being simple \(\slt\)-modules.
\end{theorem}

\begin{proof}
The critical irreducibility theorem of \cite{FGXZ} makes
\(\overline M^{\,c}_{a,b}(\chi;\boldsymbol\theta)\) simple. The \(F\)-torsion
vectors form a submodule because \(\ad F\) is locally nilpotent; the critical
PBW basis contains the nonzero vectors \(f_q^{\,m}\bar u\), and twisting back
by \(\sigma_{a+1}\) shows that \(f_b\) is not locally nilpotent on the whole
module. Hence \(f_b\) is injective. Since each \(T(m)\) commutes with \(\slh\)
and with \(F\), exactness of Ore localization gives
\[
\T_F\!\left(M^c_{a,b}/I^{a,b}_{\boldsymbol\theta}M^c_{a,b}\right)
\cong
\frac{\T_FM^c_{a,b}}{I^{a,b}_{\boldsymbol\theta}\bigl(\T_FM^c_{a,b}\bigr)}.
\]
Theorem~\ref{thm:boundary} and \(Z_i^{a,b}=Z_i^{a+1,b-1}=T(q-i)\) identify the
right side with \(\overline M^{\,c}_{a+1,b-1}(\psi;\boldsymbol\theta)\). The
\(f\)-version follows from Proposition~\ref{prop:loc-induction}, with
simplicity from \cite[Theorem 4.2]{FGXZ}.
\end{proof}

Iterating,
\[
\T_{f_{b-k+1}}\cdots\T_{f_b}
\overline M^{\,\bullet}_{a,b}(\chi;\boldsymbol\theta)
\cong
\overline M^{\,\bullet}_{a+k,b-k}(\chi^{(k)};\boldsymbol\theta),
\qquad \bullet\in\{c,f\},
\]
with \(\boldsymbol\theta\) unchanged along the entire chain.

\section{Vertex-algebra viewpoint}\label{sec:vertex}

By smoothness and the fixed action of \(K\), each \(Q_c\) is a weak module for
the universal affine vertex algebra \(V^\kappa(\slc)\); see
\cite[Section 2]{FGXZ}. Descending to its simple quotient requires a separate
check that the defining vertex ideal annihilates the module; smoothness alone
does not imply this.

\section{Open problems}\label{sec:open}

\begin{enumerate}[leftmargin=*,label=\arabic*.]
\item \textbf{Nonintegral simplicity in the basic case.}
Under \(\lambda_1\ne0\) and \(\kappa\ne-2\), the natural candidate is
\[
T_{f_0}M_{-1,1}(\chi)\text{ simple}
\quad\Longleftrightarrow\quad \lambda_0\notin\Z.
\]
Only the necessary direction is proved. The missing step is to reduce an
arbitrary nonzero submodule vector to a pure fraction that generates the
whole quotient.
\item \textbf{The corresponding deep-region criterion.}
Under \(\lambda_L\ne0\), \(\kappa\ne-2\), and \(d_*\ge L\), does
\(\mu=\lambda_0-c\kappa\notin\Z\) imply simplicity?
\item \textbf{The integral-parameter submodule structure.}
Determine whether \(G_c\) and \(Q_c/G_c\) are simple, whether \(G_c\) is the
socle, whether their extension splits, and whether \(Q_c\) has finite
Jordan--H\"older length.
\item \textbf{Simplicity in the shallow region.}
Here \(e_{-c}\) is injective for all \(\mu\), so the deep-region integral
obstruction disappears. Does simplicity of the original module suffice?
\item \textbf{Parameter isomorphisms and degenerations.}
The level \(\kappa\), support \(\lambda_0+2\Z\), negative-root
local-nilpotence set, and positive-Cartan local-finiteness profile give
necessary invariants; for \(i>d_*\), the generalized eigenvalue \(\lambda_i\)
is an additional invariant. A complete classification, especially for
\(\lambda_L=0\) or \(\kappa=-2\), remains unresolved.
\item \textbf{Critical-level central reduction.}
At \(\kappa=-2\), study the completed center, its central reductions, and
their compatibility with localization; Proposition~\ref{prop:ann} does not
settle completed-center questions.
\item \textbf{Multiple missing modes and twisted localization.}
Proposition~\ref{prop:commuting} suggests a family indexed by finite sets of
free modes. One may also study the generalized conjugation
\[
\Theta_\xi(s)=\sum_{j\ge0}\binom{\xi}{j}(\ad F)^j(s)F^{-j},\qquad \xi\in\C,
\]
on \(D_FM\), where the sum is finite for each \(s\in U\). This requires \(F\)
to be invertible; it is not directly an operation on \(Q_c\).
\end{enumerate}

\section{Sources and verification record}\label{sec:verify}

The original induced modules, their simplicity criterion, spectral-flow
automorphisms, and the smooth-module vertex-algebra correspondence are drawn
from \cite{FGXZ}; \cite{FGM} provides the localization comparison for
free-field modules. The boundary isomorphism and the derivation-induction
framework were established in the earlier project notes \cite{Boundary,
LivingV8}, and the nonboundary results in \cite{Nonboundary}.

Machine checks. The Takiff action formulas of Section~\ref{sec:basic} were
checked with 5,866 exact rational evaluations
(\texttt{work/check\_nonboundary\_takiff.py}), including the six generators'
bracket relations, the highest-vector recurrence at integral parameters,
explicit locally nilpotent vectors, and pure-fraction resonance, with
\(\lambda_1=0\) included and no polynomial-degree truncation. Independently,
in the verification round of 2026-09-02 the key formulas underlying
Sections~\ref{sec:pbw}, \ref{sec:boundary}, \ref{sec:derivation}, and
\ref{sec:critical} were checked symbolically with three engines
(Python/SymPy, Mathematica 14.3, GAP 4.15.1), 76 checks in total, all
passing: the localization formulas as a consistent Laurent action for all
\(s,t\in\Z\) (both sides are polynomials of degree at most three in \(s,t\)),
the closed commutation forms, the \(\sigma_k\) automorphism on all defining
relations, the character property of \(\psi\), the cyclic-vector relations,
the chain \(E^pF^{-1}u=(-1)^pp!\lambda_L^{\,p}F^{-p-1}u\) for a grid of
\((a,b)\), the derivation shifts, and the composition law plus \(K=-2\)
invariance of the spectral-flow formula. Two boundary facts were also
accepted into the Danus LLM-review fact graph (fact IDs
\texttt{d6bf6fef62dc8971}, \texttt{b388fabb701946d3}). No Lean/Coq
formalization is claimed, and the paper-level theorems (the FGXZ criterion,
the critical irreducibility statement, and the completed-center computations)
remain external inputs.

\begingroup
\raggedright
\begin{thebibliography}{9}
\bibitem{FGXZ}
V.~Futorny, X.~Guo, Y.~Xue, and K.~Zhao,
\emph{Smooth representations of affine Kac--Moody algebras},
arXiv:2404.03855v2 (2024; journal version 2025).
\href{https://arxiv.org/html/2404.03855v2}{arXiv full text}.

\bibitem{FGM}
V.~Futorny, D.~Grantcharov, and R.~A.~Martins,
\emph{Localization of free field realizations of affine Lie algebras},
arXiv:1404.7148v2 (2015).
\href{https://arxiv.org/html/1404.7148v2}{arXiv full text}.

\bibitem{Boundary}
Earlier project working notes, \emph{Boundary localization: advisor note}
(September 2026) and \emph{Boundary localization and derivation extension}
(September 8, 2026).

\bibitem{LivingV8}
Project living note, \emph{Smooth representations of affine Kac--Moody
algebras: PBW bases, irreducibility, and localization tasks},
version v8 (2026-09-02).

\bibitem{Nonboundary}
Project note, \emph{Nonboundary localization quotients of smooth affine
\(\mathfrak{sl}_2\)-modules} (2026-09-11).
\end{thebibliography}
\endgroup
\end{document}
